\section{Exact parameters and numerical certification}\label{sec:parameters}

We now evaluate the logarithmic yield rate $H_0+C_*$ and terminal
logarithmic volume rate $S_*$ for one choice of distributions.
The same calculation verifies positive entries
of the stage-capacity vectors, the balance conditions, and $S_*>0$, then
bounds the exponent ratio obtained from \eqref{eq:rank-constraint}.
Five integer arrays in
Appendix~\ref{app:parameter-tables}, followed by rational operations,
specify this choice; no optimization procedure is part of the certificate.

For an integer argument $n$, define the rational weight
\begin{equation}\label{eq:rational-weight}
 F(n)=\left(\sum_{j=0}^{16}\frac{(n/640000)^j}{j!}\right)^{64}.
\end{equation}
The initial weights, the two split-weight arrays, and the terminal weights
use 223 distinct arguments of $F$, including the implicit zero, all between
$-120770$ and $45052$. The binary-parameter array instead supplies the
fractions $n/10^6$ and is not passed to $F$. For the 223 weight arguments,
exact rational evaluation gives
\begin{equation}\label{eq:weight-bounds}
 2^{-23}<F(n)<2^{23}.
\end{equation}
The truncated series is positive directly as well: for nonnegative arguments
all terms are nonnegative, and for negative arguments the even Taylor partial
sum lies above the positive exponential. This observation concerns ordinary
real arithmetic used to choose probabilities, and imposes no restriction on
the field of the tensor construction.

\paragraph{Complete decoding rules.}
All shapes and pairs below are traversed in increasing lexicographic order,
and coordinate labels and array indices start at zero. For a shape $s$ of
size $2\ell$, write
\[
 \mathcal B(s)=\{b\in\mathbb Z_{\ge0}^3:
          b_0+b_1+b_2=\ell,\ b_i\le s_i\ (0\le i<3)\}.
\]
The arrays are consumed as follows.
\begin{enumerate}
\item The 17 entries $a_0,\ldots,a_{16}$ define the initial sorted law by
\[
 A_g=\frac{\operatorname{mult}(g)\prod_{i=0}^2F(a_{g_i})}
 {\sum_{h_0\le h_1\le h_2,\,\sum h_i=16}
             \operatorname{mult}(h)\prod_{i=0}^2F(a_{h_i})},
 \qquad
 \operatorname{mult}(g)=\frac{3!}{\prod_j\bigl(\#\{i:g_i=j\}\bigr)!}.
\]
Thus the law on ordered shapes, before coalescing permutations, is a product
of three one-coordinate weights.
\item The 118 entries in the Stage A array are read over all positive sorted
$g$ of size 16. The 54 entries in the Stage B array are read over all positive
\emph{ordered} $t$ of size 8. At either parent $s$, set
$\ell=\frac12\sum_i s_i$. For coordinate $i=0,1,2$, with $k=s_i$, read one
integer $n_{s,i}(j)$ for each
\[
 \max(0,k-\ell)\le j<\lfloor k/2\rfloor,
 \qquad n_{s,i}(\lfloor k/2\rfloor)=0
\]
(the final zero consumes no entry). For $b\in\mathcal B(s)$ put
\[
 w_s(b)=\prod_{i=0}^2F\bigl(n_{s,i}(\min(b_i,s_i-b_i))\bigr),
 \qquad q_s(b)=\frac{w_s(b)}{\sum_{c\in\mathcal B(s)}w_s(c)}.
\]
Use $P_g=q_g$ at Stage A and $p_t=q_t$ at Stage B. Each is a positive
product-weight law and satisfies $q_s(b)=q_s(s-b)$ exactly.
\item Read the 75 entries of the fourth array over positive ordered $t$ of
size 8, and then over $u\in\mathcal B(t)$ with $\max u=2$. An entry $n$
means $d_{t,u}=n/10^6$. These include both the zero-coordinate $022$ shapes
and the positive $112$ shapes. The range is $0\le d_{t,u}\le0.110106$.
\item Read the 36 entries of the fifth array over all ordered $t$ of size 8
with a zero coordinate. For each $t$, list the pairs
\[
 (i,j),\qquad 0\le i\le j\le5,\qquad r_i+r_j=\max t.
\]
Read an integer for each pair except the last, whose integer is zero.
Assign its $F$-weight to both $(i,j)$ and $(j,i)$ and normalize over
\emph{ordered} pairs to obtain $z_t$. Hence an off-diagonal weight occurs
twice in the denominator, while a diagonal weight occurs once. The
multiplicities $D_iD_j$ count words at the leaf and are not included in this
probability normalization.
\end{enumerate}

There are 30 sorted initial shapes, 21 positive sorted initial shapes,
21 positive ordered size-8 shapes, and 24 size-8 shapes with a zero.
All specified parents consume their entries, including parents whose eventual
amount is zero. For example, $(6,1,1)$ is included in the Stage B and binary
traversals. Skipping it, or a zero-amount parent in the $z$ traversal, would
change the meaning of the arrays. The exact types in this construction are
rational, because \eqref{eq:rational-weight} and every normalization are
rational. Only the balance fractions, and limiting tolerance choices, require
subsequent rational approximation.

The complement prescriptions in Section~\ref{sec:stagger} can be checked
without logarithms. On the six statistic slots, $\kappa$ is an involution
with $D_{\kappa(i)}=D_i$ and $r_{\kappa(i)}=4-r_i$. At a $022$ child it fixes
both slots 2 and 3, so the common parameter on the two nonzero sides gives
exactly matching laws. At zero-coordinate length-four shapes the other side
receives the pushforward of $z_t$ in both slots. These identities, including
all zero-axis cases, are checked with exact integers and fractions by the
archived verifier. Complementary \emph{child shapes} need not have equal
parameters: for instance
\[
 d_{(1,3,4),(0,2,2)}=\frac{92966}{10^6},\qquad
 d_{(1,3,4),(1,1,2)}=\frac{1204}{10^6}.
\]
They retain their separate prescriptions throughout the contractions.

\paragraph{Finite contractions.}
Here is an explicit evaluation rule for the conditional entropies appearing
in the capacities. For a finite nonnegative vector $x$, let
\[
 \mathcal K(x)=\left(\sum_i x_i\right)\log\left(\sum_i x_i\right)
                   -\sum_i x_i\log x_i,
 \qquad 0\log0=0.
\]
Thus $\mathcal K(x)=M H(x/M)$ when $M=\sum_i x_i>0$, and
$\mathcal K(0)=0$. Given a parent $s$, its split law $q$, a priority
$(X,Y,Z)$, and half-laws $Q_{b,W}$, form the ordered pair law
\[
 D_W(i,j)=\sum_{b\in\mathcal B(s)}q(b)Q_{b,W}(i)Q_{s-b,W}(j).
\]
For $W=Y$, let $\mathcal I_W=\{b:b_Z=0\}$; for $W=Z$, let
$\mathcal I_W=\{b:b_Xb_Y=0\}$. The grouped entropy is exactly
\begin{equation}\label{eq:finite-group-contraction}
 J_W=\sum_{b\in\mathcal I_W}q(b)\mathcal K(Q_{b,W})
       +\sum_{k=0}^{s_W}
        \mathcal K\left(\sum_{\substack{b\notin\mathcal I_W\\b_W=k}}
                                     q(b)Q_{b,W}\right).
\end{equation}
The capacity vector is
\[
 \left(\mathcal K\bigl((\sum_{b:b_X=k}q(b))_k\bigr),
        \mathcal K(D_Y)-2J_Y,\mathcal K(D_Z)-2J_Z\right).
\]
At A substitute $(s,q,Q)=(g,P_g,V)$ and weight by $A_g$; at B substitute
$(t,p_t,v_{t,\cdot,\cdot})$ and weight by $m_t$, with the priorities already
specified. This rule uses each parent once. Child amounts and leaf volumes
use both halves as in Proposition~\ref{prop:leaf-volume}. Along with the
formulas for $H_0,B_*,T,S_0,S_1,S_2$, it specifies only finite additions,
multiplications, divisions, and logarithms of rational numbers.

\begin{table}[htbp]
\centering
\small
\begin{tabular}{c r r r}
\hline
Partial capacity & position 0 & position 1 & position 2\\\hline
A: $\min g=1$ & .010739192 & .016912163 & .016213531\\
A: $\min g=2$ & .081941050 & .088150511 & .082591626\\
A: $\min g=3$ & .292033029 & .283922596 & .257330639\\
A: $\min g=4$ & .547143900 & .527202791 & .487211274\\
A: $\min g=5$ & .247999127 & .248080948 & .237330883\\\hline
B: $116$ & .002770751 & .002770751 & .000443544\\
B: $125$ & .080908872 & .106275766 & .087117877\\
B: $134$ & .386356324 & .454421695 & .323639129\\
B: $224$ & .395346750 & .395353092 & .248930076\\
B: $233$ & .611253765 & .686627626 & .634289551\\\hline
\end{tabular}
\caption{Partial sums defining $L_A$ and $L_B$. A B row sums over all ordered
parents with that sorted shape. Each displayed entry has absolute error less
than $4.795\cdot10^{-10}$.}
\label{tab:capacities}
\end{table}
\begin{table}[htbp]
\centering
\begin{tabular}{c r@{\qquad}c r@{\qquad}c r}
\hline
$H_0$ & 1.841147374 & $B_*$ & 1.158066972 & $T$ & 2.345955063\\
$S_0$ & .031704071 & $S_1$ & .882362236 & $S_2$ & 11.680704014\\\hline
\end{tabular}
\caption{The other scalar contractions, each with absolute error less than
$3\cdot10^{-8}$.}
\label{tab:scalars}
\end{table}

\FloatBarrier
\begin{proposition}[Numerical certificate]\label{prop:numerical-certificate}
The exact parameters above satisfy every positivity and balance condition
used in Proposition~\ref{prop:staggered-yield}. All thirty partial capacity
entries in Table~\ref{tab:capacities} are positive. In particular, $S_*>0$,
$T\log2-B_*>0$, all entries of $L_A,L_B,L_C$ are positive, and the three
balance fractions are strictly positive. Moreover,
\begin{align}
 5.612983956059236756&<H_0+C_*<5.612983956059236757,\label{eq:yield-certificate}\\
 12.594770321594676362&<S_*<12.594770321594676363,\label{eq:volume-certificate}\\
 2.371054886006745684
 &<\frac{3(8\log7-H_0-C_*)}{S_*}
       <2.371054886006745685.\label{eq:ratio-certificate}
\end{align}
With $\Omega=2.371056$, the strict comparison has the outward enclosure
\begin{equation}\label{eq:gap-certificate}
 0.000004676829725968
 < H_0+C_*+\frac{\Omega}{3}S_*-8\log7
 <0.000004676829725969.
\end{equation}
\end{proposition}

\begin{proof}
We give the elementary interval rule and its error accounting, so that the
certificate does not depend on assuming decimal logarithms are accurate.
For $1\le y\le2$ put
\[
 P(y)=2\sum_{j=0}^{30}\frac{1}{2j+1}
                    \left(\frac{y-1}{y+1}\right)^{2j+1},
 \qquad
 \mathcal L(x)=eP(2)+P(2^{-e}x),\quad e=\lfloor\log_2x\rfloor.
\]
For rational $x>0$, its exponent $e$ is determined by exact comparisons with
powers of two. Integrating the geometric series for $2/(1-z^2)$ gives
\[
 0\le\log y-P(y)\le
 \tau:=\frac{2\,3^{-63}}{63(1-1/9)}<3.120348\cdot10^{-32}.
\]
Thus $|\log x-\mathcal L(x)|\le(|e|+1)\tau$, including when $e<0$.
On each dyadic interval $\mathcal L$ is increasing; at a boundary its left
limit is $eP(2)+P(2)=(e+1)P(2)$, the value from the next interval.
It is therefore continuous and increasing on $(0,\infty)$. Applying $P$
at outward endpoints, then enlarging by the indicated tail, encloses
$\log x$ even if an input interval crosses a dyadic boundary.

For completeness, all exponents in this evaluation satisfy $|e|<2048$.
There are fewer than 256 ordered points in every product-law support.
By \eqref{eq:weight-bounds}, a nonzero product-law probability exceeds
$2^{-146}$, a $z$ entry exceeds $2^{-54}$, and a positive binary entry
exceeds $2^{-20}$. The initial marginal entries exceed $2^{-148}$.
Positive entries of a length-four law $V$ exceed $2^{-186}$; multiplying
one by a Stage A probability gives $2^{-332}$, and multiplying two gives
$2^{-518}$. These bounds also cover sums within nonzero groups; all these
masses are at most one. The integer leaf multiplicities and coefficients
are less than $7^8<2^{23}$. Hence the actual logarithm arguments lie between
$2^{-518}$ and $2^{23}$, a stronger range than needed. One may uniformly use
\begin{equation}\label{eq:log-error}
 \epsilon=2048\tau<6.391\cdot10^{-29}.
\end{equation}

Here is an independent propagation bound for replacing every logarithm
literally by $\mathcal L$, without presuming that the resulting rational
version of $\mathcal K$ is homogeneous. For a vector of total mass $M$, the
error in $\mathcal K$ is at most $2M\epsilon$. The total singleton and
nonsingleton group masses in \eqref{eq:finite-group-contraction} sum to one,
so its error is at most $2\epsilon$. A sharing capacity has errors at most
$(2,6,6)\epsilon$ in its three positions. The A parent amounts sum to at
most one, the B parent amounts to at most two, and the positive length-two
amount $T$ is at most four. Consequently the entrywise error in $L_A+L_B$
is at most $(6,18,18)\epsilon$, and its average has error at most
$14\epsilon$. The error in $h_b(d)$ is at most $3\epsilon$, hence that in
$B_*$ at most $20\epsilon/3$. Including $H_0$ gives
\[
 |\text{error in }(H_0+C_*)|\le\tfrac{68}{3}\epsilon.
\]
For $S_0,S_1,S_2$ the corresponding bounds are
$\epsilon,6\epsilon,20\epsilon$, respectively. Therefore
\[
 |\text{error in }S_*|\le27\epsilon,\qquad
 |\text{error in the comparison at }\Omega|
 \le\left(\tfrac{68}{3}+9\Omega+8\right)\epsilon
 <3.324\cdot10^{-27}.
\]

The file \texttt{verification/scripts/verify\_parameters.py} implements the
stated finite contractions. It constructs every $F(n)$ as an exact fraction,
then uses outward interval operations for all arithmetic, including each
normalizing denominator. Shared occurrences of a denominator are retained
as interval occurrences; no independence or cancellation is assumed to
reduce error. Exact structural checks establish normalization by division
by the positive total, complement symmetry, and statistic support before
the numerical checks. The interval checks of zero differences are only
consistency checks of these identities.
The primary verification evaluates the rational series above and adds its
tail; a second run uses directed native logarithms as a cross-check.
Binary interval endpoints are converted to exact fractions and serialized
by floor for lower bounds and ceiling for upper bounds, including negative
endpoints. These operations give outward decimal bounds without a
nearest-rounding assumption. At 80 decimal digits, every reported native-log
interval has width below $10^{-76}$; with the rational-series tail included,
every reported interval has width below $3\cdot10^{-29}$. These bounds include
all normalization and subsequent arithmetic rounding.

In addition to \eqref{eq:yield-certificate}--\eqref{eq:gap-certificate},
the resulting bounds include
\[
 0.468025165562828972<T\log2-B_*<0.468025165562828973.
\]
The following entries enclose $L_C$, the numerators
$b_i=T\log2-C_*+(L_A+L_B)_i$ of the balance fractions, and the fractions
$\lambda_i$, respectively. Each listed number is a lower bound; adding
$10^{-15}$ to it gives an upper bound.
\[
\begin{array}{c|rrr}
 &0&1&2\\\hline
 L_C&1.115343821954349&.962118642424118&1.396738451843254\\
 b&.510748315682386&.663973495212618&.229353685793481\\
 \lambda&.363761291246081&.472890166361139&.163348542392779
\end{array}
\]
The fractions sum to one exactly by their definition, since the sum of
their numerators is $3(T\log2-B_*)$. Positivity of $L_A,L_B$ follows
already from the thirty positive partial entries. The boundary deficit
used in the finite-lot schedule also satisfies
\[
 0.533798153<2C_*-\min_i(L_A)_i-\min_i(L_A+L_B)_i
 -\min_i(L_B+L_C)_i-\min_i(L_C)_i<0.533798155<1.
\]
These bounds verify all stated feasibility conditions. The script checks
every displayed enclosure and stores the full interval outputs
with the packaged integer data.

The rounded entries in Tables~\ref{tab:capacities} and~\ref{tab:scalars}
already certify $\omega<2.371056$.  The stronger bound in
Theorem~\ref{thm:main} uses the interval enclosure
\eqref{eq:ratio-certificate}.
For the rounded-table check, allow independent errors of $4\cdot10^{-8}$ in all 30 partial capacities, $H_0$, $B_*$, and the three
volume rows. The total uncertainty in the comparison is at most
\[
 \left(12+\Omega\right)4\cdot10^{-8}=5.7484224\cdot10^{-7}.
\]
Using the displayed centers and an outward bound on $8\log7$ leaves a
positive margin greater than $4.1011\cdot10^{-6}$. Thus the strict inequality
has ample room for both arithmetic enclosure and the arbitrarily small
losses in the construction.
\end{proof}

\begin{proof}[Proof of Theorem~\ref{thm:main}]
Since $S_*>0$, Proposition~\ref{prop:numerical-certificate} and
\eqref{eq:rank-constraint} give
\[
 \omega\le w\le\frac{3(8\log7-H_0-C_*)}{S_*}
 <2.371054886006746<2.371054887.\qedhere
\]
\end{proof}
The integer tables and their full traversal are sufficient to reproduce this
comparison. The verification package includes a readable assignment table and both
interval outputs; verification requires Python 3.10 or later and
\texttt{mpmath} 1.3.0, with no optimization software or external data.
